% Shared preamble for the two point-by-point response letters (one per reviewer).
\documentclass[11pt,a4paper]{article}
\usepackage[margin=2.4cm]{geometry}
\usepackage{amsmath,amssymb}
\usepackage[T1]{fontenc}
\usepackage{lmodern}
\usepackage{xcolor}
\usepackage{framed}
\usepackage{enumitem}
\usepackage[hidelinks]{hyperref}
\DeclareMathOperator{\spec}{spec}
\DeclareMathOperator{\diag}{diag}
\setlength{\parskip}{4pt}
\setlength{\parindent}{0pt}
\definecolor{shadecolor}{gray}{0.93}
% \comment{label}{text}: the reviewer's comment, quoted in a shaded box.
\newcommand{\comment}[2]{\par\medskip\begin{shaded}\textbf{Comment #1.} \emph{#2}\end{shaded}}
% \response{text}
\newcommand{\response}[1]{\par\textbf{Response.} #1\par\medskip}
% \where{text}: pointer into the revised manuscript (section, equation, page/line numbers of the revised PDF).
\newcommand{\where}[1]{\par\textit{Location in the revised manuscript:} #1\par}
\newcommand{\manuscriptid}{mathematics-4633263}
\newcommand{\manuscripttitle}{From Cain to Matignon: Exact Positive-Diagonal Stability in Dimension Three and a Nonlinear Caputo Ecological Realization}
\newcommand{\letterhead}[1]{%
  \begin{center}
  {\Large\bfseries Response to #1}\\[6pt]
  Manuscript \texttt{\manuscriptid}, \emph{Mathematics} (MDPI)\\[2pt]
  \emph{\manuscripttitle}
  \end{center}
  \medskip
  Page and line numbers below refer to the revised manuscript (line numbers printed in the margin of the revised PDF). Every change requested by #1 is addressed individually; changes made in response to the other reviewer are listed at the end for completeness, because they affect the same sections.\par\medskip}

\begin{document}
\letterhead{Reviewer 2}

We thank the reviewer for the assessment and for the seven points, each of which is addressed below. Points 3 and 4 led to new text and, for point 4, to new certified and numerical results; point 6 led to a substantially extended numerical study; the remaining points are editorial.

\comment{1}{On page 5, line 177, ``Therefore, the fractional-only'' instead of ``Therefore the fractional-only''.}
\response{Corrected.}
\where{Section 3, proof of Theorem 1, page 5, line 179 of the revised manuscript.}

\comment{2}{On page 22, line 683, ``consistent with the positive real part'' instead of ``consistently with the positive real part''.}
\response{Corrected.}
\where{Section 7.3, page 25, line 739.}

\comment{3}{On page 8, line 245, the main characterization in Theorem 2 is established under the assumption that \(-A\) is a strict \(P\)-matrix. The authors should explain the implications of this restriction and clarify whether the characterization can be extended to matrices with zero principal minors.}
\response{We added an explicit discussion in two places. First, the opening of Section 4.3 states that the strict-\(P\) hypothesis isolates the open, full-dimensional principal-minor stratum to which Theorem 2 applies, and that the theorem does not classify matrices with vanishing signed principal minors. Second, a new paragraph after Lemma 2 (principal-minor necessity) answers the extension question in three steps.
\begin{enumerate}[leftmargin=1.6em]
\item \emph{Implication for the main result.} Matrices with a vanishing signed principal minor lie in the algebraic hypersurfaces \(\{a_{ii}=0\}\) and \(\{\det A[I]=0\}\). These strata are closed, have Lebesgue measure zero, and by Lemma 2 contain no interior point of \(\mathcal F_\alpha^{(3)}\). The characterization of the full-dimensional interior, which is the object of the paper, is therefore complete as stated.
\item \emph{The strata are not trivial.} One cannot say that nothing happens on them. Embedding the two-dimensional fractional-only class of Theorem 1 as
\[
A=\begin{pmatrix}0&b&0\\c&0&0\\0&0&-1\end{pmatrix},\qquad bc<0,
\]
gives \(\spec(DA)=\{\pm i\sqrt{-d_1d_2bc},\,-d_3\}\) for every positive diagonal \(D\). Hence \(A\in\mathcal F_\alpha^{(3)}\setminus\mathcal D_H^{(3)}\) for every \(0<\alpha<1\), although \(p_1=p_2=0\): a fractional-only matrix on the excluded stratum.
\item \emph{Why the characterization does not extend directly.} The four-invariant method degenerates there: if \(p_i=0\), the simplex coordinate \(x_i=p_id_i/a_D\) vanishes identically and the invariants \(\beta_I=m_I/\prod_{i\in I}p_i\) are undefined, while \(m_I=0\) removes the corresponding monomial from the orbit coefficients. A classification of these lower-dimensional strata would require separate face-by-face limiting arguments. The revised text states that this is not attempted in the present paper.
\end{enumerate}
We did not add numerical experiments for this point, because a numerical scan can neither establish nor rule out an extension of an exact all-diagonal characterization.}
\where{Section 4.3, page 9, lines 270--275 (scope statement) and the paragraph ``Vanishing principal minors'', page 9, lines 286--297.}

\comment{4}{The manuscript establishes local spectral stability for the fractional-order ecological system. The authors should discuss whether the positive equilibrium remains biologically feasible under perturbations in the model parameters and initial conditions.}
\response{We now treat the two kinds of perturbation separately, and in both cases we replaced a qualitative statement by an explicit one.

\emph{Parameters.} At the realization W2-A the coexistence equilibrium satisfies \(X,Y,Z>0\) and \(m<X<K\), the scalar equilibrium equation has a simple root, and both invariant inequalities have strict margins, so the implicit-function theorem preserves a nearby feasible equilibrium and the strict band under sufficiently small perturbations of the 14 biological parameters. To make ``sufficiently small'' explicit, we certified a neighborhood with the same outward-rounded interval arithmetic used for the \(m\)-branch: all 14 parameters are varied independently in \([p_i(1-\delta),p_i(1+\delta)]\) around the values of Table 1 (with \(e_1\) and \(e_3\) varied separately), at \(\alpha=0.9\). For \(\delta=0.005\) (199 subdivision boxes) and for \(\delta=0.0075\) (28\,802 boxes), every box is certified to contain a unique prey root with \(X,Y,Z>0\), \(m<X<K\) and \(s>0\), hence a strict-\(P\) reduced matrix, and to satisfy both strict inequalities \(\kappa-T_1(\beta)>0\) and \(T_{0.9}(\beta)-\kappa>0\). The resulting enclosures are \(X\in[0.981,1.021]\), \(Y\in[0.756,0.813]\) and \(Z\in[0.170,0.212]\). Thus every parameter vector within \(0.75\%\) relative distance of W2-A in each coordinate has a feasible positive coexistence equilibrium in the open fractional-only band. The previously certified segment \(0.394\le m\le0.578\) in the single direction \(m\) remains in the text.

\emph{Initial conditions.} The submitted manuscript only said that small positive perturbations remain positive ``locally in time'' by continuous dependence; this was weaker than what holds, and we replaced it. The vector field is \(C^1\) near \(u^*\), and at every parameter vector of the certified box the Jacobian \(RJ=R\diag(u^*)B\) is Matignon stable for every positive diagonal response-rate matrix \(R\), in particular for \(R=I\) and for the configuration \(R_*\) used in Section 7.3. The linearized asymptotic stability theorem for Caputo systems (Cong, Doan, Siegmund and Tuan, \emph{Electron. J. Qual. Theory Differ. Equ.} 2016, No.~39, Theorem 3.1; new reference [27]) therefore applies: for every \(\varepsilon>0\) there is \(\delta>0\) such that \(\|u(0)-u^*\|<\delta\) implies \(\|u(t)-u^*\|\le\varepsilon\) for all \(t\ge0\) and \(u(t)\to u^*\); choosing \(\varepsilon<\min_iu_i^*\) keeps all three densities positive for all time. The text states that the theorem does not quantify \(\delta\). As numerical evidence only, explicitly not as a proof, we added an integration of 1024 random relative perturbations of \(u^*\) at each of 20 radii between \(0.01\) and \(0.9\) at \(\alpha=0.9\), with the same PECE scheme and step refinement as in Section 7.3 and horizon \(t=600\). With the configuration \(R_*\) every trajectory remained positive and returned to the equilibrium for relative radius up to \(0.35\) (at radius \(0.40\), 11 of 1024 trajectories did not return); with the unscaled model \(R=I\) the corresponding radius was \(0.15\) (87 failures at \(0.20\)). Every trajectory that did not return approached the extinction state, as expected from the strong Allee threshold. The text states that these radii are finite-sample, finite-horizon estimates that depend on the response-rate configuration.

The Discussion was updated accordingly. We added no new figure or table for this point, and the local and spectral scope of the main result is restated: no global positivity or permanence conclusion is claimed.}
\where{Section 7, page 22, lines 633--637 (feasibility and implicit-function statement); Section 7.1, paragraph ``Certified 14-parameter neighborhood'', page 23, lines 667--678; Section 7.3, paragraph ``Initial-condition perturbations'', page 25, lines 718--736; Section 8, page 29, lines 806--815.}

\comment{5}{Figure captions should be sufficiently descriptive to allow the figures to be understood independently.}
\response{All captions were rewritten so that each figure can be read on its own. Each caption now identifies the plotted objects and axes, names the mathematical boundaries shown (Cain boundary \(\kappa=T_1\), Matignon boundary \(\kappa=T_\alpha\), feasibility edge), states which markers are interval-certified and which are numerical evaluations of exact formulas, and states the limitation of finite numerical search where relevant. The artwork of Figures 1 and 2 was regenerated to remove internal labels, and the new Figure 7 has a self-contained caption.}
\where{Captions of Figures 1--8 (pages 6, 11, 17, 21, 24, 27, 27 and 28).}

\comment{6}{Section 7 presents a nonlinear simulation for \(\alpha=0.9\) and compares it with the classical case \(\alpha=1\). The authors already report time-step refinement results, which is useful. Nevertheless, simulations at additional fractional orders, particularly near the predicted stability boundary, would provide stronger numerical support for the theoretical conclusions.}
\response{We added a full nonlinear Caputo study across the predicted boundary order. For the selected positive response-rate configuration \(R_*\), the theorem predicts loss of Matignon stability exactly at the order \(\alpha_*\) with \(T_{\alpha_*}(\beta)=\kappa\), and \(\alpha_*=0.9505501071\ldots\) coincides to the displayed precision with \((2/\pi)\min_j|\arg\lambda_j(D_*B)|\). The new study integrates the same nonlinear model, with the same response rates and the same initial perturbation, for \(0.7\le\alpha\le1\): a dense grid with \(\Delta\alpha=0.001\) and additional points around \(\alpha_*\), 363 integrations in total, each with the Adams--Bashforth--Moulton PECE scheme at \(h=0.02\) and \(h=0.01\) for step refinement, horizon \(t=600\) away from the boundary and \(t=2000\) for \(|\alpha-\alpha_*|<0.02\), and termination when the relative distance from the equilibrium first reaches the declared local radius \(0.2\). An independent adaptive integration at \(\alpha=1\) reproduces the integer-order exit time.

Results: the fitted late log-distance rate changes sign between \(\alpha=0.9500\) (negative) and \(\alpha=0.9506\) (positive), and every order \(\alpha\ge0.953\) exits the local radius before \(t=2000\) (at \(t=1955\) for \(\alpha=0.953\), \(t=1084\) for \(0.955\), \(t=101.6\) for \(\alpha=1\)). The end-versus-start distance diagnostic agrees with the linear spectral prediction in 339 of 340 dense runs; the single exception, \(\alpha=0.9506\), has a positive late rate but has not yet recovered its initial distance by \(t=2000\), and the text uses it to warn against classifying by final-versus-initial distance alone. The unscaled configuration \(R=I\), whose Jacobian is Hurwitz, decays at every tested order. The new Table 2 lists representative orders and the new Figure 7 shows the trajectory envelopes on both sides of the boundary and the fitted rate as a function of \(\alpha\), with the analytic threshold drawn as a vertical line that is not fitted to the trajectories. The text states that this is finite-horizon numerical corroboration of the spectral classification and not an independent proof of asymptotic stability. A preliminary short-horizon table (\(t=180\)) prepared during the revision was not retained, because final distances alone did not distinguish slower decay from instability near the boundary.}
\where{Section 7.3, paragraph ``Long-horizon nonlinear Caputo verification across the exact boundary order'', page 25, lines 742 ff.; Table 2 (page 26); Figure 7 (page 27); Section 8, page 27.}

\comment{7}{Many references are very old. Add a few recent and relevant studies on the stability of nonlinear fractional-order systems and clearly highlight the novelty of the proposed approach.}
\response{We added recent, published, DOI-verified studies on stability and nonlinear dynamics of fractional-order ecological systems, including commensurate and incommensurate predator--prey models with Allee effects, fear effects, refuge, delays and ecoepidemiological coupling (references [11--14] in the Introduction: Liu, Li and Huang, \emph{Math.\ Comput.\ Simul.} 2025; Alomari and Bahaa, \emph{Sci.\ Rep.} 2026; Baghel et al., \emph{Chaos Solitons Fractals} 2026; Saha et al., \emph{Fract.\ Calc.\ Appl.\ Anal.} 2026), together with recent double-Allee predator--prey references in Section 6 (\emph{Fractal Fract.} 2026, \emph{J.\ Biol.\ Phys.} 2025, \emph{AIMS Math.} 2026, \emph{Qual.\ Theory Dyn.\ Syst.} 2025) and the 2016 linearized-stability theorem used for point 4. The Introduction now states the novelty sharply: these studies concern specific models or prescribed parameters and do not yield an exact all-positive-diagonal spectral classification for arbitrary strict-\(P\) real three-by-three matrices, whereas the present paper eliminates the universal positive-diagonal quantifier exactly into four invariants and one scalar threshold that deforms Cain's classical criterion, and then realizes the resulting open band by a nonlinear Caputo ecological model. We kept the classical references (Cain 1976, Kellogg 1972, Matignon 1996) because the results build on them directly.}
\where{Section 1, page 3, lines 82--91; reference list, pages 29--30.}

\section*{Changes made in response to Reviewer 1 that affect the same sections}
For the reviewer's information: the proofs of Lemma 1, Corollary 1 and Theorem 5 and the attainment of the minimum defining \(T_\alpha\) were made explicit (Sections 4.1, 4.2, 4.3 and 5.2); the exact fractional-only set \(T_1\le\kappa<T_\alpha\) was distinguished from its interior (Section 4.3); the unscaled Jacobian \(J=\diag(u^*)B\) was distinguished from the row-scaled Jacobian \(R_*J=D_*B\) used in the time-domain comparison (Sections 7.2--7.3); the interval-certification procedure, working precisions and the validated lower bound of the threshold were described in Section 7.1, and the printed value of \(m_0\) was aligned with the stated precision; internal labels were removed from the text and the figures. None of these changes alters a theorem statement.
\end{document}
